\documentclass[12pt]{article}

\usepackage[a4paper, total={6in, 8in}]{geometry}
\usepackage{textcomp}

\begin{document}
\setlength{\parindent}{0pt}

\def \t {\textrightarrow}
\def \v {\vspace{1em}}
\def \bi {\begin{itemize}}
\def \ei {\end{itemize}}
\def \s[#1] {\section*{#1}}
\def \ss[#1] {\subsection*{#1}}
\def \sss[#1] {\subsubsection*{#1}}

\s[Italo Svevo]
Zeno Cosini è l'alter ego di svevo \\
Duplice natura, zeno viene da xenus: straniero, mentre cosini: piccola cosa \t mostra già dal nome l'abbassamento tonale che si contrappone a d'annuznio \\
Il nome è il presagio dell'anti eroe \t questo personaggio è l'immagine dello svevo autore \t simbiosi tra i due \\
Questa immagine ci viene proiettata attraverso il contesto del periodo storico che presagisce il conflitto mondiale \\
Il narratore e l autore vivono una simbiosi di critica e criticità \t attraverso presentazione anche ironica \\
La coscienza non è salvata \t il personaggio chiude senza aver preso consapevolezza di cio che lo ha portato a intraprendere terapia \\
C'è un presupposto pero: la terapia per curare una malattia psichica \t ci si domanda pero che cosa la società consegna a noi di un "sano", o di un "malato" \\
Il malato e il sano scardinano la concezione d'annunziana \t i piu sani sono in realta i veri malati \t il concetto di salute assume diverse sfumature in diverse epoche \\
La salute interiore non sempre coincide con la salute esteriore, vitalismo \t i veri sani sono i finti malati di fatto, ovvero coloro che hanno consapevolezza della loro salute precaria \\
Il focus sulle figure femminili si snodano nel orizzonte di zeno, seguono questo concetto di salute \t le donne sane incarnano i valori borghesi del tempo \t l'etichetta prospetta la sanita \\
Assenza di alcune malattie trasmissive \t perbenismo che caratterizza l'epoca, che porta alle scelte \\
Il personaggio non sa scegliere, ma si lascia scegliere \t lo stesso zeno cosini, e lo stesso svevo (che lavora nella ditta del suocero, ma non sceglie di lavorare qua, lascia scegliere ad altri \t si voleva dedicare solo alla letteratura)

\v

Questo da il la al romanzo \t segue il criterio della libera associazione di pensieri (non ti tutte le terapie) \t mostra la precarieta della psicanalisi \\
La libera associazione si mette in campo quando il paziente parla a ruota libera \t il terapeuta pero non da risposte, ma induce noi stessi a trovarle \\
Poderosi balzi indietro, e poi fa ripiombare nel presente \t e poi ci sono le valutazioni del terapeuta, zeno pian piano fa riaffiorare ricordi ma non sempre in modo autentico \\
Gli inganni ricordano anche il mondo pirandelliano \t il personaggio alla fine annienta la persona \t rivoluzione copernicana delle opere di pirandello \\
Pird. pero aveva salvato la coscienza, qua invece viene subissata \t i finti sani sono i veri malati \t la prestanza di alcuni è la vera malattia \\
Sono malati del credersi sani, di non porsi domande, di accettare le ottiche borghesi che deresponsabilizzano \\
La reale identita interiore è quella che alla fine permette di far coincidere quello che è esposto con quello che esiste? \t l'autenticita viene compromessa \\
La vicenda si apre sul fumo \t in questo evento \t lui è un terapeuta, non uno psichiatra \t e lo psichiatra è anche diverso anche dallo psicologo

\v

Il fatto di non riuscire a smettere di fumare \t ???? \t l'inconscio vuole emergere, con tutti gli atti mancati e gli irrisolti \\
L'effetto pigmalione è il copione che si realizza, il copione dell'inconscio \\
L'incoscio porta al replicarsi di di alcune azioni che non ci spieghiamo consciamente \\
Il copione che si autoavvera \\
Qua la pscianalisi non è scienza, mentre al giorno d'oggi si \\
Sono strutture interiori queste, non cristallizzazione della forma \\
A trieste sono presenti gli allievi di freud \t il cognato era curato dal dottor Weiss \t e svevo lo accompagnava alle sedute \\
S. è sigmund probabilmente \t continuo rimbalzo tra il Weiss, e freud stesso \t jung è stato anche lui suo allievo, ma ha preso percorso diverso \\
Cambiano i tempi, e anche l'approccio diventa anche in parte modificato \t anche nell'approcciarsi al paziente \\
Dopo aver identificato nel padre la ragione della sua dipendenza al fumo, capisce quanto erano fallaci i suoi rapporti amorosi \\
Anche in una borghesia apparentemente sana, perde il mordente \\
Questo porta alla scelta che chiude l'opera di rinunciare l'inettitudine \t zeno è un inetto, che però si sente vincente \t trova il coraggio di dire no a una terapia che non gli va bene \\
In questo caso reagisce, pero mostrando una volonta di non curarsi \t l umanita è inquinata alle radici, quando tutto verra disintegrato solo allora l uomo si potra ricostruire \\
Il male è in tutte le cose, non è salvezza per l uomo \t bisogna struggere tutto, anche le piante

\v

Si snoda da questo il tema della guerra \t finisce la terapia con la scusa della guerra \t ma in realta lui non lo vuole pagare \\
Idea del remunerare investe tutte le discipline che trattano anime \\
Questo svilisce molto la figura di zeno cosini, che appare come un uomo che non vuole mostrare il vero lato di se \\
Zeno mostra cio che non è, la volonta di conoscersi \t ma all epilogo si mostra la non volonta di conoscere se stesso \\
Il terapeuta sollicita domande che non vogliamo porsi, non da risposte \\
Noi ci accomodiamo alla fine nell'abitus che all inizio volevamo toglierci \\
Inabilita a vivere, vivere è infatti decidere di fare qualcosa \\
Lui non sceglie, ma si lascia decidere \t è la morte dell'identita dell'essere
\end{document}
